\section{From controlled fillings to finite skeleta}
\label{sec:products}

Proposition~\ref{prop:uniform-fillings} concerns ordinary integral
chains.  We will apply it coordinate by coordinate to products of
regular modules, keeping the finite set of simplex types common to
all coordinates.  We first explain why the resulting homology
vanishing yields geometric, rather than merely homological,
finiteness.

All coefficient modules in this section are right modules, and all
group homology is integral.  For a group $J$ and a set $\mathcal I$,
write $(\Z J)^{\mathcal I}$ for the direct \emph{product} of copies
of the right regular module, with coordinatewise action.

\subsection{A geometric product-module criterion}

The following argument combines the product-module finiteness method of
Bieri--Eckmann and Brown~\cite{BieriEckmann1974,Brown1975} with the
cellular passage to geometric finiteness.  We give the argument in the
form needed here, including its finite-presentation hypothesis.

\begin{lemma}\label{lem:product-criterion}
Let $J$ be a finitely presented group.  If
\[
 H_n\bigl(J,(\Z J)^{\mathcal I}\bigr)=0
 \qquad\text{for every set }\mathcal I\text{ and every }n\ge2,
\]
then $J$ is of type $\Finfty$.
\end{lemma}

\begin{proof}
We construct nested, $n$-dimensional CW complexes $X_n$, for
$n\ge2$, with fundamental group $J$, finite $n$-skeleton, and
$(n-1)$-connected universal cover.  Each passage from $X_n$ to
$X_{n+1}$ will attach finitely many $(n+1)$-cells.  A finite
presentation complex supplies $X_2$.

Suppose that $X_n$ has been constructed, and put
$Z=\widetilde X_n$.  Thus $Z$ is an $n$-dimensional,
$(n-1)$-connected free $J$-complex.  By successively attaching
$J$-orbits of cells of dimensions at least $n+1$ to kill its
homotopy groups, extend $Z$ to a contractible free $J$-complex $Y$.
There is no finiteness requirement on these auxiliary attachments.
Put $R=\Z J$.  The augmented cellular chain complex $C_*(Y)$ is a
free left $R$-resolution of $\Z$, and $C_q(Y)=C_q(Z)$ is a
finite-rank free $R$-module for $q\le n$.

Let $\mathcal I$ be the set of all cellular $n$-cycles of $Z$, and
let $M=R^{\mathcal I}$.  If $C$ is a finite-rank free left
$R$-module, the coordinate map is an isomorphism
\begin{equation}\label{eq:finite-rank-product-tensor}
 M\otimes_R C\longrightarrow C^{\mathcal I},
 \qquad (m_z)_{z\in\mathcal I}\otimes c
       \longmapsto(m_zc)_{z\in\mathcal I}.
\end{equation}
Using this identification in degrees $n$ and $n-1$, the family
whose $z$-coordinate is the cycle $z$ defines a cycle
$\xi\in M\otimes_R C_n(Y)$.  The hypothesis $H_n(J,M)=0$ gives
\[
 \xi=\partial y
 \qquad\text{for some }y\in M\otimes_R C_{n+1}(Y).
\]
Choose a free $R$-basis of $C_{n+1}(Y)$.  A tensor is a finite sum,
and each chain uses only finitely many basis elements.  Consequently
we can write
\[
 y=\sum_{j=1}^s m_j\otimes c_j
\]
using finitely many basis elements $c_1,\ldots,c_s$, even though
$C_{n+1}(Y)$ need not have finite rank.  Taking the $z$-coordinate
of the boundary equation gives
\[
 z=\sum_{j=1}^s (m_j)_z\,\partial c_j.
\]
Thus the module of cellular $n$-cycles in $Z$ is generated over
$R$ by $\partial c_1,\ldots,\partial c_s$.

Because $Z$ is $(n-1)$-connected and $n\ge2$, the natural Hurewicz map
$\pi_n(Z)\to H_n(Z;\Z)$ is an isomorphism
\cite[Theorem~4.37 and p.~370]{Hatcher2002}.  Naturality with respect to
deck transformations makes it an isomorphism of $R$-modules, using the
canonical change of basepoint in the simply connected space $Z$.
Since $Z$ has no $(n+1)$-cells, its $n$th homology is its module of
cellular $n$-cycles.  Choose maps of $n$-spheres representing the displayed
finite list of generators, and attach their $J$-translates to $Z$.
Downstairs this attaches finitely many $(n+1)$-cells to $X_n$;
upstairs the resulting complex is $n$-connected.  It is the
universal cover of the resulting complex $X_{n+1}$, since cells of
dimension at least three do not change the fundamental group.
This completes the induction.

Let $X=\bigcup_{n\ge2}X_n$.  It has fundamental group $J$ and
finitely many cells in every dimension.  The universal cover of
$X$ is weakly contractible.  It is connected and simply connected, and
for $q\ge2$, cellular approximation \cite[Theorem~4.8]{Hatcher2002}
homotopes any map $S^q\to\widetilde X$ into the $q$-skeleton of
$\widetilde X$, which is already present in $\widetilde X_q$.  That
class is killed in
$\widetilde X_{q+1}$.  A weakly
contractible CW complex is contractible by Whitehead's theorem
\cite[Theorem~4.5]{Hatcher2002}, so $X$ is a $K(J,1)$.
\end{proof}

\subsection{The ascending HNN sequence}

For the rest of the section assume the hypotheses of
Theorem~\ref{thm:torus}, and set $J=T(U,f)$.  We will compute the
homology in Lemma~\ref{lem:product-criterion} using a tree for $J$.
We give the construction explicitly to fix the direction of the
coefficient map.

Form the direct limit
\[
 L=\varinjlim\bigl(U\xrightarrow{f}U\xrightarrow{f}U
          \xrightarrow{f}\cdots\bigr),
\]
and denote its stage maps by $\iota_i:U\to L$, for $i\ge0$.
They satisfy $\iota_i=\iota_{i+1}f$.  Since $f$ is injective,
each $\iota_i$ is injective and
$L=\bigcup_{i\ge0}\iota_i(U)$ is an increasing union.  The map
\[
 \sigma\bigl(\iota_i(u)\bigr)=\iota_{i+1}(u)
\]
is an automorphism: its inverse sends $\iota_i(u)$ to
$\iota_i(f(u))$.  In the semidirect product adjoining a letter
$t$ with $t x t^{-1}=\sigma(x)$ for $x\in L$, we have
\[
 t^{-1}\iota_0(u)t=\iota_0(f(u)).
\]
This gives a homomorphism from $J$ to that semidirect product.
Conversely, the maps $\iota_i(u)\mapsto t^iut^{-i}$ are compatible
with the direct system and with $\sigma$, and give the inverse
homomorphism.  Thus $U$ embeds in $J$.  The same description shows
that every element of $J$ can be written
\begin{equation}\label{eq:ascending-normal-form}
 t^p u t^{-q},\qquad u\in U,\quad p,q\ge0.
\end{equation}
Indeed, an element of $L\rtimes\langle t\rangle$ has the form
$\iota_p(u)t^m$; increasing $p$ if necessary makes $q=p-m$
nonnegative.

Consider the graph whose vertices and oriented edges are both
indexed by the left cosets $jU$, with the edge indexed by $jU$
running from $jU$ to $jtU$.  The terminal vertex is well defined
because
\begin{equation}\label{eq:ascending-commutation}
 ut=tf(u)\qquad(u\in U).
\end{equation}
The homomorphism $J\to\Z$ that kills $U$ and sends $t$ to $1$
defines a height on the vertices.  Each edge increases height by
one, and each vertex has exactly one outgoing edge.
If $j=t^p u t^{-q}$ is as in
\eqref{eq:ascending-normal-form}, the forward path of length $q$
from $jU$ reaches $t^pU$.  Every forward ray therefore meets the
ray $U,tU,t^2U,\ldots$, proving connectedness.  An undirected
cycle would have a vertex of least height with two outgoing edges,
which is impossible.  The graph is a tree.  Left multiplication
gives a $J$-action with one vertex orbit and one edge orbit, and
both the chosen vertex $U$ and the chosen edge from $U$ to $tU$
have stabilizer $U$.

Let $C_*(U)$ be the normalized homogeneous resolution introduced
in the preceding section.  Since $\Z J$ is free as a right
$\Z U$-module, induction preserves its exactness.  We obtain the free left
$\Z J$-resolution
\[
 P_* =\Z J\otimes_{\Z U}C_*(U)
\]
of the permutation module $\Z[J/U]$.  The two endpoint maps on
this resolution are the identity and the map
\[
 E(j\otimes c)=jt\otimes f_*(c).
\]
The second is well defined: for $u\in U$,
\[
 jut\otimes f_*(c)
   =jt f(u)\otimes f_*(c)
   =jt\otimes f_*(uc).
\]
It is a left $\Z J$-linear chain map, lifting the terminal-vertex
map $jU\mapsto jtU$.

The mapping cone of $1-E:P_*\to P_*$, with its natural
augmentation, is a free left $\Z J$-resolution of $\Z$.
To see exactness, first take the homology of $P_*$: it is
concentrated in degree zero and equal to $\Z[J/U]$.  The
remaining two-term complex is the augmented cellular chain
complex of the tree, which is exact.  Equivalently, the long exact
sequence of the cone identifies its homology with that of this
two-term complex.

For an arbitrary right $\Z J$-module $M$, tensoring this free
resolution with $M$ yields the long exact sequence
\begin{equation}\label{eq:ascending-hnn-homology}
 \begin{split}
 H_n(U,M)&\xrightarrow{\,1-F\,}H_n(U,M)
      \longrightarrow H_n(J,M)\\
      &\longrightarrow H_{n-1}(U,M)
          \xrightarrow{\,1-F\,}H_{n-1}(U,M),
 \end{split}
\end{equation}
where $M$ is restricted to $U$ in the outer terms.  On the chain
complex $M\otimes_{\Z U}C_*(U)$, the map $F$ is
\begin{equation}\label{eq:coefficient-shift}
 F(m\otimes c)=mt\otimes f_*(c).
\end{equation}
In particular the coefficient multiplier is $t$, not $t^{-1}$.

\subsection{Products of regular modules}

Fix any set $\mathcal I$ and take $M=(\Z J)^{\mathcal I}$.
We claim that $F$ induces a locally nilpotent endomorphism of
$H_i(U,M)$ for every $i>0$: each individual homology class is
killed by some power of $F$.

Recall that a bound is a finite symmetric subset $A\subseteq U$
containing $1$, and $C_q^A(U)$ is generated by the normalized
tuples $(u_0,\ldots,u_q)$ with $u_i^{-1}u_j\in A$ for all $i,j$.
Let $\mathcal S_q(A)$ be their finite set of types up to left
translation.  The first vertex can be translated to $1$, after
which all other vertices lie in $A$, proving finiteness.  As
$C_q^A(U)$ is free over $\Z U$ on these types, there are natural
identifications
\begin{equation}\label{eq:bounded-product-chains}
 \begin{split}
 M\otimes_{\Z U}C_q^A(U)
 &\cong\bigoplus_{s\in\mathcal S_q(A)}M\\
 &\cong\prod_{\alpha\in\mathcal I}
      \left(\bigoplus_{s\in\mathcal S_q(A)}\Z J\right).
 \end{split}
\end{equation}
The second equality interchanges a product with a \emph{finite}
direct sum; no analogous interchange with an infinite direct sum
is being used.

Here is a concrete interpretation of
\eqref{eq:bounded-product-chains}.  The free abelian basis of
$\Z J\otimes_{\Z U}C_q^A(U)$ consists of normalized tuples with
$u_0,\ldots,u_q\in U$ of the form
\[
 (ju_0,\ldots,ju_q),\qquad j\in J,
 \quad u_a^{-1}u_b\in A\text{ for all }a,b.
\]
Each tuple lies in one left coset of $U$.  Expanding the group-ring
coefficients, a tensor chain therefore gives, in each coordinate
$\alpha$, a finite ordinary integral chain in such coset tuples.
There is one bound $A$ for every coordinate.  Conversely, any
family of coordinatewise finite chains with a common bound
defines a tensor chain by \eqref{eq:bounded-product-chains}.
Their lengths, coefficients, and group-ring supports need not
admit any bound independent of $\alpha$.

Every chain in $M\otimes_{\Z U}C_q(U)$ uses finitely many
simplex types and consequently belongs to a bounded subcomplex
of this form.  Indeed, take the union of the finitely many
relative vertex differences in its tensor factors, then add
their inverses and $1$.  The inclusion of this bounded complex
is injective degree by degree, since its types are a subset of
the free basis of $C_q(U)$.

Now let $z\in M\otimes_{\Z U}C_i(U)$ be a cycle, where $i>0$,
and choose one input bound $A$ for it.  In each coordinate
$\alpha$, split its ordinary chain into its finitely many
left-coset components.  Each component is a cycle, since the
boundary preserves the coset and the coset chain groups form a
direct sum.  Choose an origin $j$ in each nonzero component's
coset.  Left translation by $j^{-1}$ identifies the component
with an ordinary cycle in $C_i^A(U)$.

Proposition~\ref{prop:uniform-fillings} supplies an exponent
$k\ge0$ and a finite output bound $B$ that work for
\emph{every} such cycle, in every coordinate and every coset.
By \eqref{eq:coefficient-shift}, the map $F^k$ changes the
origin $j$ to $jt^k$ and applies $f^k$ to the relative vertices.
Indeed, on tuples the single-step formula is simply
\[
 (ju_0,\ldots,ju_i)
    \longmapsto(jt f(u_0),\ldots,jt f(u_i))
    =(ju_0t,\ldots,ju_it).
\]
This also verifies that the formula is independent of the
chosen origin.

For each component, choose an ordinary finite filling of its
relative image under $f^k$ in $C_{i+1}^B(U)$, and translate
this filling by $jt^k$.  Sum the fillings within each coordinate.
The result is still finite in that coordinate because there were
only finitely many nonzero coset components.  Different input
cosets may have the same output coset; in that case their
fillings are simply added.  Across all coordinates the bound
is still the single finite set $B$.
Equation~\eqref{eq:bounded-product-chains} therefore assembles
these fillings into a chain
\[
 y\in M\otimes_{\Z U}C_{i+1}^B(U).
\]
The boundary and $F^kz$ lie in the bounded complex for
$D=B\cup f^k(A)$, a finite symmetric set containing $1$.  They agree
in every coordinate by construction.  The coordinate identification
\eqref{eq:bounded-product-chains} for $D$ is injective, so
$\partial y=F^kz$ in the tensor complex.
Thus $F^k[z]=0$, proving the claimed local nilpotence.  This
argument uses ordinary fillings chosen separately in each coset;
it requires no equivariance of the homotopies or fillings in
Proposition~\ref{prop:uniform-fillings}.

\subsection{Completion of the torus theorem}

If an endomorphism $F$ of an abelian group is locally nilpotent,
then $1-F$ is invertible.  On an element $x$, its inverse is
\[
 x+Fx+F^2x+\cdots,
\]
where only finitely many terms are nonzero.  The sum defines a
homomorphism, and telescoping verifies both inverse identities.
It follows that the two maps $1-F$ adjacent to $H_n(J,M)$ in
\eqref{eq:ascending-hnn-homology} are isomorphisms for $n\ge2$.
Exactness gives
\[
 H_n\bigl(J,(\Z J)^{\mathcal I}\bigr)=0
 \qquad(n\ge2).
\]
Only positive-degree local nilpotence was needed: no assertion
about $H_0(U,M)$ enters this conclusion.

By Lemma~\ref{lem:torus-presentation}, $J$ is finitely presented.
The set $\mathcal I$ was arbitrary, so
Lemma~\ref{lem:product-criterion} proves that $J$ is of type
$\Finfty$.  The direct-limit construction already proved that
$U$ embeds in $J$.  This completes the proof of
Theorem~\ref{thm:torus}.
